\subsection{Tensor Product Bases and Maps}

\begin{definition} \label{def:finite-coordinate-tensor-product-model} \hlt{Finite-Coordinate Tensor Product Model}\\
Let $\beta=(\mathbf{v}_1,\ldots,\mathbf{v}_m)$ and $\gamma=(\mathbf{w}_1,\ldots,\mathbf{w}_n)$ be ordered bases of $\mathcal{V}$ and $\mathcal{W}$.\\
The \hlt{coordinate tensor model} is $\mathcal{V}\otimes_{\beta,\gamma}\mathcal{W}\defeq\operatorname{Mat}(m,n,\mathbb{F})$. For $(\mathbf{v},\mathbf{w})\in\mathcal{V}\times\mathcal{W}$, define $\mathbf{v}\otimes_{\beta,\gamma}\mathbf{w}\defeq[\mathbf{v}]_\beta[\mathbf{w}]_\gamma^{\mathsf{T}}$.\\
The identity $\mathbf{v}_i\otimes_{\beta,\gamma}\mathbf{w}_j=E_{i,j}$ holds by the coordinate columns of the basis vectors.
\end{definition}

\begin{proposition} \label{prop:tensor-product-matrix-basis-model} \hlt{Tensor Model Basis and Bilinear Evaluation}\\
In the coordinate tensor model associated to bases $\beta=(\mathbf{v}_1,\ldots,\mathbf{v}_m)$ and $\gamma=(\mathbf{w}_1,\ldots,\mathbf{w}_n)$:
\begin{enumerate}[label=(\roman*)]
\setlength{\itemsep}{0pt}
\item The matrices $E_{i,j}=\mathbf{v}_i\otimes_{\beta,\gamma}\mathbf{w}_j$ form a basis.
\item The map $(\mathbf{v},\mathbf{w})\mapsto\mathbf{v}\otimes_{\beta,\gamma}\mathbf{w}$ is bilinear.
\item Let $\mathcal{X}$ be $\mathbb{F}$-vector space and let $f\in\operatorname{Bil}_{\mathbb{F}}(\mathcal{V},\mathcal{W};\mathcal{X})$.\\
A unique $\Theta_f\in\LL(\operatorname{Mat}(m,n,\mathbb{F}),\mathcal{X})$ satisfies $(\forall\,(\mathbf{v},\mathbf{w})\in\mathcal{V}\times\mathcal{W})\,\Theta_f(\mathbf{v}\otimes_{\beta,\gamma}\mathbf{w})=f(\mathbf{v},\mathbf{w})$.
\end{enumerate}
\begin{proof}
\begin{enumerate}[label=(\roman*)]
\setlength{\itemsep}{0pt}
\item The matrices $E_{i,j}$ form the entry-matrix basis of $\operatorname{Mat}(m,n,\mathbb{F})$ by Proposition \ref{prop:matrix-spaces-are-vector-spaces}(ii).
\item Fix $a,b\in\mathbb{F}$, $\mathbf{v},\mathbf{v}'\in\mathcal{V}$, and $\mathbf{w},\mathbf{w}'\in\mathcal{W}$. The identity $(a\mathbf{v}+b\mathbf{v}')\otimes_{\beta,\gamma}\mathbf{w}=a(\mathbf{v}\otimes_{\beta,\gamma}\mathbf{w})+b(\mathbf{v}'\otimes_{\beta,\gamma}\mathbf{w})$ holds by $\mathbb{F}$-linearity of $[\cdot]_\beta$.\\
Likewise, $\mathbf{v}\otimes_{\beta,\gamma}(a\mathbf{w}+b\mathbf{w}')=a(\mathbf{v}\otimes_{\beta,\gamma}\mathbf{w})+b(\mathbf{v}\otimes_{\beta,\gamma}\mathbf{w}')$ by $\mathbb{F}$-linearity of $[\cdot]_\gamma$.
\item The basis-value assignment $E_{i,j}\mapsto f(\mathbf{v}_i,\mathbf{w}_j)$ extends uniquely to $\Theta_f\in\LL(\operatorname{Mat}(m,n,\mathbb{F}),\mathcal{X})$ by Proposition \ref{prop:linear-map-determined-by-basis}.\\
Write $\mathbf{v}=\sum\limits_{i=1}^{m}x_i\mathbf{v}_i$ and $\mathbf{w}=\sum\limits_{j=1}^{n}y_j\mathbf{w}_j$. The identity $\Theta_f(\mathbf{v}\otimes_{\beta,\gamma}\mathbf{w})=\sum\limits_{i=1}^{m}\sum\limits_{j=1}^{n}x_iy_jf(\mathbf{v}_i,\mathbf{w}_j)$ holds by linearity of $\Theta_f$.\\
Last sum equals $f(\mathbf{v},\mathbf{w})$ by bilinearity of $f$. Uniqueness holds by agreement on the basis matrices $E_{i,j}$ and Proposition \ref{prop:linear-map-determined-by-basis}.
\end{enumerate}
\hlt{Structural Commentary}: Entry matrices form the basis; coordinate expansion proves factorisation.\\
\hlt{Key Insight}: Each basis pair occupies one independent matrix coordinate.
\end{proof}
\end{proposition}

\begin{illustrationfigure}[H]
  \centering
  \begin{tikzpicture}[
  >=Stealth,
  every node/.style={font=\small},
  map/.style={->, gray!65!black, line width=0.75pt},
  slot/.style={draw=gray!45, fill=gray!4, minimum width=0.68cm, minimum height=0.52cm},
  active/.style={slot, fill=teal!10, draw=teal!65!black, line width=0.65pt},
  basis/.style={draw=gray!45, fill=gray!3, rounded corners=3pt, line width=0.55pt,
    minimum width=2.45cm, minimum height=0.78cm, align=center}
]
  \node[basis] (pair) at (-2.45,0.52) {$\mathbf{v}_i\otimes_{\beta,\gamma}\mathbf{w}_j$};

  \draw[map] (pair.east) -- (-0.60,0.52);

  \foreach \x in {0,1,2,3} {
    \foreach \y in {0,1,2} {
      \node[slot] at (\x*0.70,\y*0.54) {};
    }
  }
  \node[active] at (1.40,0.54) {$E_{i,j}$};

  \draw[gray!50] (-0.02,0.54) -- (2.10,0.54);
  \draw[gray!50] (1.40,-0.26) -- (1.40,1.34);
\end{tikzpicture}

  \caption{A basis pair occupies the $(i,j)$ matrix slot.}
  \label{fig:tensor-product-matrix-basis}
\end{illustrationfigure}

Figure \ref{fig:tensor-product-matrix-basis} illustrates the basis correspondence in Proposition \ref{prop:tensor-product-matrix-basis-model}(i). The proposition owns the basis and bilinearity claims.

\begin{theorem} \label{thm:tensor-bases-dimension-and-rank} \hlt{Tensor Bases, Dimension, and Coordinate Rank}\\
Let $\mathcal{U}$ and $\mathcal{V}$ be finite-dimensional with bases $\beta=(\mathbf{u}_1,\ldots,\mathbf{u}_m)$ and $\gamma=(\mathbf{v}_1,\ldots,\mathbf{v}_n)$.
\begin{enumerate}[label=(\roman*)]
\setlength{\itemsep}{0pt}
\item The tensors $\mathbf{u}_i\otimes\mathbf{v}_j$ form a basis of $\mathcal{U}\otimes\mathcal{V}$.
\item $\dim(\mathcal{U}\otimes\mathcal{V})=\dim(\mathcal{U})\dim(\mathcal{V})$.
\item Every $\mathbf{z}\in\mathcal{U}\otimes\mathcal{V}$ admits unique coefficients $r_{ij}\in\mathbb{F}$. They satisfy $\mathbf{z}=\sum\limits_{i=1}^{m}\sum\limits_{j=1}^{n}r_{ij}(\mathbf{u}_i\otimes\mathbf{v}_j)$.
\item Define $R_{\beta,\gamma}(\mathbf{z})\defeq[r_{ij}]$ from (iii). The value $\operatorname{rank}R_{\beta,\gamma}(\mathbf{z})$ is independent of the bases $\beta,\gamma$.
\end{enumerate}
\begin{proof}
\begin{enumerate}[label=(\roman*)]
\setlength{\itemsep}{0pt}
\item The coordinate tensor model is a tensor product by Proposition \ref{prop:tensor-product-matrix-basis-model}(i)--(iii) and Definition \ref{def:tensor-product-universal-property}(i).\\
A unique comparison isomorphism $\Xi:\mathcal{U}\otimes\mathcal{V}\to\operatorname{Mat}(m,n,\mathbb{F})$ exists by Corollary \ref{cor:tensor-product-unique-up-to-isomorphism}.\\
It satisfies $\Xi(\mathbf{u}_i\otimes\mathbf{v}_j)=E_{i,j}$ by the same corollary. The tensors $\mathbf{u}_i\otimes\mathbf{v}_j$ form a basis by bijectivity of $\Xi$ and Proposition \ref{prop:tensor-product-matrix-basis-model}(i).
\item $\dim(\mathcal{U}\otimes\mathcal{V})=mn=\dim(\mathcal{U})\dim(\mathcal{V})$ by the $mn$-element basis in (i).
\item The coefficients $r_{ij}$ exist uniquely by the basis in (i).
\item Let $\beta',\gamma'$ be other ordered bases and set $P\defeq[I_{\mathcal{U}}]_{\beta'\leftarrow\beta}$ and $Q\defeq[I_{\mathcal{V}}]_{\gamma'\leftarrow\gamma}$.\\
The identity $R_{\beta',\gamma'}(\mathbf{z})=P R_{\beta,\gamma}(\mathbf{z})Q^{\mathsf{T}}$ holds by factorwise coordinate substitution.\\
Matrices $P,Q^{\mathsf{T}}$ invertible by Theorem \ref{thm:change-of-basis-matrices}(i), Corollary \ref{cor:transpose-invertible-iff}. Equality $\operatorname{rank}R_{\beta',\gamma'}(\mathbf{z})=\operatorname{rank}R_{\beta,\gamma}(\mathbf{z})$ holds by Theorem \ref{thm:full-rank-products-preserve-rank}(ii).
\end{enumerate}
\hlt{Structural Commentary}: Coordinates transfer the basis; invertible basis changes preserve rank.\\
\hlt{Key Insight}: Basis changes multiply coordinate matrices by invertible factors.
\end{proof}
\end{theorem}

\begin{definition} \label{def:tensor-coordinate-matrix-and-rank} \hlt{Tensor Coordinate Matrix and Rank}\\
Let $\mathcal{U},\mathcal{V}$ be finite-dimensional, with bases $\beta=(\mathbf{u}_1,\ldots,\mathbf{u}_m)$ and $\gamma=(\mathbf{v}_1,\ldots,\mathbf{v}_n)$.
\begin{enumerate}[label=(\roman*)]
\setlength{\itemsep}{0pt}
\item For $\mathbf{z}\in\mathcal{U}\otimes\mathcal{V}$, let $r_{ij}$ be the unique coefficients from Theorem \ref{thm:tensor-bases-dimension-and-rank}(iii). The \hlt{coordinate matrix} is $R_{\beta,\gamma}(\mathbf{z})\defeq[r_{ij}]$.
\item The \hlt{rank of the tensor} is $\operatorname{rank}(\mathbf{z})\defeq\operatorname{rank}R_{\beta,\gamma}(\mathbf{z})$. It is independent of $\beta,\gamma$ by Theorem \ref{thm:tensor-bases-dimension-and-rank}(iv).
\end{enumerate}
\end{definition}

\begin{theorem} \label{thm:tensor-zero-test-and-decomposable-rank} \hlt{Zero Tests and Decomposable Rank}\\
Let $\mathcal{U},\mathcal{V}$ be $\mathbb{F}$-vector spaces.
\begin{enumerate}[label=(\roman*)]
\setlength{\itemsep}{0pt}
\item $(\forall\,\mathbf{u}\in\mathcal{U})(\forall\,\mathbf{v}\in\mathcal{V})\,\mathbf{0}\otimes\mathbf{v}=\mathbf{u}\otimes\mathbf{0}=\mathbf{0}$.
\item Assume Zorn's lemma, let $r\in\mathbb{Z}_{>0}$, and let $(\mathbf{u}_i)_{i=1}^{r}\in\mathcal{U}^r$ be linearly independent.\\
Let $(\mathbf{v}_i)_{i=1}^{r}\in\mathcal{V}^r$. If $\sum\limits_{i=1}^{r}\mathbf{u}_i\otimes\mathbf{v}_i=\mathbf{0}$, then $\mathbf{v}_1=\cdots=\mathbf{v}_r=\mathbf{0}$.
\item If $\dim\mathcal{U},\dim\mathcal{V}<\infty$ and $\mathbf{z}\in(\mathcal{U}\otimes\mathcal{V})\setminus\{\mathbf{0}\}$, then $\mathbf{z}$ is decomposable $\Longleftrightarrow\operatorname{rank}(\mathbf{z})=1$.
\item Assume $\dim\mathcal{U},\dim\mathcal{V}<\infty$ and let $\mathbf{z}\in\mathcal{U}\otimes\mathcal{V}$. For $r\in\mathbb{Z}_{\geq0}$, define $\mathcal{D}_r\defeq\{\sum\limits_{i=1}^{r}\mathbf{x}_i\otimes\mathbf{y}_i:(\mathbf{x}_i,\mathbf{y}_i)_{i=1}^{r}\in(\mathcal{U}\times\mathcal{V})^r\}$.\\
Then $\operatorname{rank}(\mathbf{z})=\min\limits\{r\in\mathbb{Z}_{\geq0}:\mathbf{z}\in\mathcal{D}_r\}$.
\end{enumerate}
\begin{proof}
\begin{enumerate}[label=(\roman*)]
\setlength{\itemsep}{0pt}
\item[$(i)$] The tensor map is linear in each slot by Definition \ref{def:tensor-product-universal-property}(i). The equality $\mathbf{0}\otimes\mathbf{v}=\mathbf{u}\otimes\mathbf{0}=\mathbf{0}$ holds by Proposition \ref{prop:linear-maps-send-zero-to-zero}.
\item[$(ii)$] Fix $k\in\{1,\ldots,r\}$ and set $\mathcal{S}_k\defeq\operatorname{span}\{\mathbf{u}_i:i\neq k\}$. The inclusion $\mathbf{u}_k\notin\mathcal{S}_k$ holds by linear independence.\\
A functional $\varphi_k\in\mathcal{U}'$ satisfies $\varphi_k|_{\mathcal{S}_k}=0$ and $\varphi_k(\mathbf{u}_k)=1$ by Proposition \ref{prop:functionals-separate-points}(iii).\\
The bilinear map $(\mathbf{u},\mathbf{v})\mapsto\varphi_k(\mathbf{u})\mathbf{v}$ induces $T_k\in\LL(\mathcal{U}\otimes\mathcal{V},\mathcal{V})$ by Theorem \ref{thm:tensor-linearizes-bilinear-maps}(ii).\\
Then $\mathbf{0}=T_k(\sum\limits_{i=1}^{r}\mathbf{u}_i\otimes\mathbf{v}_i)=\sum\limits_{i=1}^{r}\delta_{ki}\mathbf{v}_i=\mathbf{v}_k$. The equality $\mathbf{v}_1=\cdots=\mathbf{v}_r=\mathbf{0}$ holds by the arbitrary choice of $k$.
\item[$(iii)$] ($\Rightarrow$) Fix ordered bases $\beta,\gamma$ and write $\mathbf{z}=\mathbf{x}\otimes\mathbf{y}\neq\mathbf{0}$. The vectors $\mathbf{x},\mathbf{y}$ are nonzero by (i).\\
The identity $R_{\beta,\gamma}(\mathbf{z})=[\mathbf{x}]_\beta[\mathbf{y}]_\gamma^{\mathsf{T}}$ holds by Definition \ref{def:tensor-coordinate-matrix-and-rank}(i) and bilinearity.\\
The equality $\operatorname{rank}(\mathbf{z})=1$ holds by Definition \ref{def:tensor-coordinate-matrix-and-rank}(ii) and $\operatorname{col}(R_{\beta,\gamma}(\mathbf{z}))=\operatorname{span}([\mathbf{x}]_\beta)\neq\{\mathbf{0}\}$.
\item[$(iii)$] ($\Leftarrow$) Fix ordered bases $\beta,\gamma$ and assume $\operatorname{rank}R_{\beta,\gamma}(\mathbf{z})=1$. Choose a nonzero column $\mathbf{a}$ spanning $\operatorname{col}(R_{\beta,\gamma}(\mathbf{z}))$.\\
A row $\mathbf{b}^{\mathsf{T}}$ satisfies $R_{\beta,\gamma}(\mathbf{z})=\mathbf{a}\mathbf{b}^{\mathsf{T}}$ by Theorem \ref{thm:full-rank-factorization}(i). Let $[\mathbf{x}]_\beta=\mathbf{a}$ and $[\mathbf{y}]_\gamma=\mathbf{b}$. Then $\mathbf{z}=\mathbf{x}\otimes\mathbf{y}$ by Theorem \ref{thm:tensor-bases-dimension-and-rank}(iii).
\item[$(iv)$] Fix ordered bases $\beta,\gamma$, set $R\defeq R_{\beta,\gamma}(\mathbf{z})$,  $r\defeq\operatorname{rank}(R)$. If $r=0$, then $\operatorname{col}(R)=\{\mathbf{0}\}$ by Definition \ref{def:column-space-matrix-rank-nullity}(iii), Proposition \ref{prop:dimension-bounds-and-basis-criteria}(ii).\\
Then $R=0$ by $\operatorname{col}(R)=\{\mathbf{0}\}$. Equalities $\mathbf{z}=\mathbf{0}$, $\mathbf{z}\in\mathcal{D}_0$ hold by Theorem \ref{thm:tensor-bases-dimension-and-rank}(iii), empty-sum convention respectively.\\
Suppose $r\geq1$, and let the columns of $X\in\operatorname{Mat}(\dim\mathcal{U},r,\mathbb{F})$ form a basis of $\operatorname{col}(R)$.\\
A matrix $Y$ satisfies $R=XY$ by Theorem \ref{thm:full-rank-factorization}(i).\\
If $\mathbf{a}_i$ is column $i$ of $X$ and $\mathbf{b}_i^{\mathsf{T}}$ is row $i$ of $Y$, then $R=\sum\limits_{i=1}^{r}\mathbf{a}_i\mathbf{b}_i^{\mathsf{T}}$ by matrix multiplication.\\
Choose $[\mathbf{x}_i]_\beta=\mathbf{a}_i$ and $[\mathbf{y}_i]_\gamma=\mathbf{b}_i$. Then $\mathbf{z}=\sum\limits_{i=1}^{r}\mathbf{x}_i\otimes\mathbf{y}_i$ by Theorem \ref{thm:tensor-bases-dimension-and-rank}(iii).\\
Membership $\mathbf{z}\in\mathcal{D}_r$ holds by definition. If $\mathbf{z}\in\mathcal{D}_s$, choose $\mathbf{z}=\sum\limits_{i=1}^{s}\mathbf{x}_i\otimes\mathbf{y}_i$. Every column of $R$ lies in $\operatorname{span}\{[\mathbf{x}_1]_\beta,\ldots,[\mathbf{x}_s]_\beta\}$.\\
The inequality $r\leq s$ holds by Definition \ref{def:column-space-matrix-rank-nullity}(iii), Proposition \ref{prop:dimension-bounds-and-basis-criteria}(i), and Theorem \ref{thm:finitedimlenofvec}.\\
The minimum formula holds by $\mathbf{z}\in\mathcal{D}_r$ and $(\forall\,s\in\mathbb{Z}_{\geq0})\,[\mathbf{z}\in\mathcal{D}_s\Rightarrow r\leq s]$.
\end{enumerate}
\hlt{Structural Commentary}: Linearity proves zeros; functionals isolate terms; matrix factors fix rank.\\
\hlt{Key Insight}: Tensor rank is the minimum number of decomposable summands.
\end{proof}
\end{theorem}

\begin{theorem} \label{thm:tensor-product-of-linear-maps} \hlt{Tensor Product of Linear Maps}\\
Let $A\in\LL(\mathcal{U},\mathcal{U}')$ and $B\in\LL(\mathcal{V},\mathcal{V}')$.
\begin{enumerate}[label=(\roman*)]
\setlength{\itemsep}{0pt}
\item A unique $A\otimes B\in\LL(\mathcal{U}\otimes\mathcal{V},\mathcal{U}'\otimes\mathcal{V}')$ exists. It satisfies $(\forall\,\mathbf{u}\in\mathcal{U})(\forall\,\mathbf{v}\in\mathcal{V})\,(A\otimes B)(\mathbf{u}\otimes\mathbf{v})=A\mathbf{u}\otimes B\mathbf{v}$.
\item Define $\tau:\LL(\mathcal{U},\mathcal{U}')\times\LL(\mathcal{V},\mathcal{V}')\to\LL(\mathcal{U}\otimes\mathcal{V},\mathcal{U}'\otimes\mathcal{V}')$ by $\tau(A,B)\defeq A\otimes B$. The map $\tau$ is bilinear.
\item Let $A_i\in\LL(\mathcal{U}_{i-1},\mathcal{U}_i)$ and $B_i\in\LL(\mathcal{V}_{i-1},\mathcal{V}_i)$ for $i\in\{1,2\}$. Then $(A_2A_1)\otimes(B_2B_1)=(A_2\otimes B_2)(A_1\otimes B_1)$.
\item $I_{\mathcal{U}}\otimes I_{\mathcal{V}}=I_{\mathcal{U}\otimes\mathcal{V}}$.
\item Suppose $\mathcal{U},\mathcal{V},\mathcal{U}',\mathcal{V}'$ are finite-dimensional with ordered bases $\beta,\gamma,\beta',\gamma'$, respectively.\\
For the lexicographically ordered tensor-product bases, $[A\otimes B]_{\beta'\otimes\gamma'\leftarrow\beta\otimes\gamma}=[A]_{\beta'\leftarrow\beta}\otimes[B]_{\gamma'\leftarrow\gamma}$.\\
The right-hand $\otimes$ denotes the matrix Kronecker product.
\end{enumerate}
\begin{proof}
\begin{enumerate}[label=(\roman*)]
\setlength{\itemsep}{0pt}
\item The map $(\mathbf{u},\mathbf{v})\mapsto A\mathbf{u}\otimes B\mathbf{v}$ is bilinear by linearity of $A,B$ and bilinearity of the tensor map.\\
The unique linear map $A\otimes B$ exists by Theorem \ref{thm:tensor-linearizes-bilinear-maps}(ii).
\item Fix $a,b\in\mathbb{F}$, $A,C\in\LL(\mathcal{U},\mathcal{U}')$, and $B,D\in\LL(\mathcal{V},\mathcal{V}')$.\\
On decomposable tensors, $\tau(aA+bC,B)=a\tau(A,B)+b\tau(C,B)$ by the linearity of $A,C$ and the tensor map.\\
Likewise, $\tau(A,aB+bD)=a\tau(A,B)+b\tau(A,D)$ on decomposable tensors.\\
The identities extend to $\mathcal{U}\otimes\mathcal{V}$ by Theorem \ref{thm:tensor-product-quotient-construction}(i), and $\tau$ is bilinear by these identities.
\item For $(\mathbf{u},\mathbf{v})\in\mathcal{U}_0\times\mathcal{V}_0$, both sides send $\mathbf{u}\otimes\mathbf{v}$ to $A_2A_1\mathbf{u}\otimes B_2B_1\mathbf{v}$.\\
The composition identity holds by decomposable spanning in Theorem \ref{thm:tensor-product-quotient-construction}(i).
\item Both $I_{\mathcal{U}}\otimes I_{\mathcal{V}}$ and $I_{\mathcal{U}\otimes\mathcal{V}}$ fix every decomposable tensor. Their equality holds by decomposable spanning in Theorem \ref{thm:tensor-product-quotient-construction}(i).
\item Fix expansions $A\mathbf{u}_j=\sum\limits_i a_{ij}\mathbf{u}'_i$ and $B\mathbf{v}_\ell=\sum\limits_k b_{k\ell}\mathbf{v}'_k$. Then $(A\otimes B)(\mathbf{u}_j\otimes\mathbf{v}_{\ell})=\sum\limits_{i,k}a_{ij}b_{k\ell}(\mathbf{u}'_i\otimes\mathbf{v}'_k)$ by bilinearity.\\
The claimed matrix identity holds by this Kronecker block-entry rule and Definition \ref{def:matrix-representation-of-linear-map}.
\end{enumerate}
\hlt{Structural Commentary}: Linearization builds $A\otimes B$; spanning proves composition; bases fix the matrix.\\
\hlt{Key Insight}: Factorwise action becomes Kronecker multiplication in coordinates.
\end{proof}
\end{theorem}

\begin{theorem} \label{thm:tensor-product-of-map-spaces} \hlt{Tensor Product of Spaces of Linear Maps}\\
Assume Zorn's lemma, and let $\mathcal{U},\mathcal{V},\mathcal{X},\mathcal{Y}$ be $\mathbb{F}$-vector spaces.
\begin{enumerate}[label=(\roman*)]
\setlength{\itemsep}{0pt}
\item A unique linear map $\Theta:\LL(\mathcal{U},\mathcal{X})\otimes\LL(\mathcal{V},\mathcal{Y})\to\LL(\mathcal{U}\otimes\mathcal{V},\mathcal{X}\otimes\mathcal{Y})$ exists. It is characterized by $\Theta(A\otimes B)(\mathbf{u}\otimes\mathbf{v})=A\mathbf{u}\otimes B\mathbf{v}$.
\item The map $\Theta$ is injective.
\item If $\mathcal{U},\mathcal{V},\mathcal{X},\mathcal{Y}$ are finite-dimensional, then $\Theta$ is an isomorphism.
\item Define $\Delta(\varphi\otimes\psi)(\mathbf{u}\otimes\mathbf{v})\defeq\varphi(\mathbf{u})\psi(\mathbf{v})$. This rule determines an injective map $\Delta:\mathcal{U}'\otimes\mathcal{V}'\to(\mathcal{U}\otimes\mathcal{V})'$.
\item If $\mathcal{U},\mathcal{V}$ are finite-dimensional, then $\Delta$ is an isomorphism.
\end{enumerate}
\begin{proof}
\begin{enumerate}[label=(\roman*)]
\setlength{\itemsep}{0pt}
\item The assignment $(A,B)\mapsto A\otimes B$ is bilinear by Theorem \ref{thm:tensor-product-of-linear-maps}(ii).\\
The unique linear map $\Theta$ exists by Theorem \ref{thm:tensor-linearizes-bilinear-maps}(ii). Its characterization holds by Theorem \ref{thm:tensor-product-of-linear-maps}(i).
\item Fix $0\neq Z=\sum\limits_{i=1}^{r}A_i\otimes B_i$ with $r$ minimal. The relations $(\forall\,i\in\{1,\ldots,r\})\,[A_i\neq0\land B_i\neq0]$ hold by minimality.\\
If $(A_i)_{i=1}^{r}$ is dependent, reindex with $A_r=\sum\limits_{i=1}^{r-1}c_iA_i$. Then $Z=\sum\limits_{i=1}^{r-1}A_i\otimes(B_i+c_iB_r)$, contrary to minimality.\\
The family $(A_i)_{i=1}^{r}$ is independent by contradiction. If $(B_i)_{i=1}^{r}$ is dependent, reindex with $B_r=\sum\limits_{i=1}^{r-1}d_iB_i$.\\
Then $Z=\sum\limits_{i=1}^{r-1}(A_i+d_iA_r)\otimes B_i$, contrary to minimality. The family $(B_i)_{i=1}^{r}$ is independent by contradiction.\\
Choose $\mathbf{u}\in\mathcal{U}$ with $A_1\mathbf{u}\neq0$. After reindexing, let $A_1\mathbf{u},\ldots,A_m\mathbf{u}$ be a basis of $\operatorname{span}\{A_i\mathbf{u}:1\leq i\leq r\}$.\\
For $m<j\leq r$, write $A_j\mathbf{u}=\sum\limits_{i=1}^{m}c_{ij}A_i\mathbf{u}$. Set $C_i\defeq B_i+\sum\limits_{j=m+1}^{r}c_{ij}B_j$ for $1\leq i\leq m$. Then $C_i\neq0$ by independence of $(B_i)_{i=1}^{r}$.\\
For $\mathbf{v}\in\mathcal{V}$, $\Theta(Z)(\mathbf{u}\otimes\mathbf{v})=\sum\limits_{i=1}^{m}(A_i\mathbf{u})\otimes C_i\mathbf{v}$. If $\Theta(Z)=0$, then $(\forall\,i\in\{1,\ldots,m\})(\forall\,\mathbf{v}\in\mathcal{V})\,C_i\mathbf{v}=\mathbf{0}$ by Theorem \ref{thm:tensor-zero-test-and-decomposable-rank}(ii).\\
The conclusion $(\forall\,i\in\{1,\ldots,m\})\,C_i=0$ holds by Assumption \ref{assump:function-extensionality}, contrary to the preceding nonvanishing.\\
The inequality $\Theta(Z)\neq0$ holds by contradiction, and $\Theta$ is injective.
\item Both spaces have dimension $\dim\mathcal{U}\dim\mathcal{V}\dim\mathcal{X}\dim\mathcal{Y}$ by Theorem \ref{thm:tensor-bases-dimension-and-rank}(ii) and Theorem \ref{thm:linear-maps-and-matrices-same-data}(ii).\\
The map $\Theta$ is an isomorphism by (ii) and Corollary \ref{cor:equal-dimension-kernel-range}(ii).
\item Set $\mathcal{X}=\mathcal{Y}=\mathbb{F}$. The multiplication map induces $\lambda:\mathbb{F}\otimes\mathbb{F}\to\mathbb{F}$ with $\lambda(a\otimes b)=ab$ by Theorem \ref{thm:tensor-linearizes-bilinear-maps}(ii).\\
The map $\lambda$ is an isomorphism by the basis correspondence $\{1\otimes1\}\mapsto\{1\}$ and Theorem \ref{thm:tensor-bases-dimension-and-rank}(i).\\
Postcomposition with $\lambda$ is an isomorphism $\lambda_\ast:\LL(\mathcal{U}\otimes\mathcal{V},\mathbb{F}\otimes\mathbb{F})\to(\mathcal{U}\otimes\mathcal{V})'$. Its inverse is postcomposition with $\lambda^{-1}$.\\
The identity $\Delta=\lambda_\ast\circ\Theta$ holds on decomposable tensors by the defining formulas. The map $\Delta$ is injective by (ii).
\item If $\mathcal{U},\mathcal{V}$ are finite-dimensional, then $\Theta$ is an isomorphism by (iii).\\
The map $\Delta$ is an isomorphism by $\Delta=\lambda_\ast\circ\Theta$ and the isomorphisms $\lambda_\ast,\Theta$.
\end{enumerate}
\hlt{Structural Commentary}: Minimality proves injection; dimensions and multiplication prove isomorphism.\\
\hlt{Key Insight}: Independent operator factors are detected on one decomposable tensor.
\end{proof}
\end{theorem}

\begin{proposition} \label{prop:scalar-extension-by-tensor-product} \hlt{Scalar Extension by Tensor Product}\\
Let $K/\mathbb{F}$ be a field extension, regard $K$ as an $\mathbb{F}$-vector space, and let $\mathcal{V}_{\mathbb{F}}$ be an $\mathbb{F}$-vector space.
\begin{enumerate}[label=(\roman*)]
\setlength{\itemsep}{0pt}
\item The $\mathbb{F}$-vector space $K\otimes_{\mathbb{F}}\mathcal{V}_{\mathbb{F}}$ is a $K$-vector space via $\alpha\cdot(\beta\otimes\mathbf{v})\defeq(\alpha\beta)\otimes\mathbf{v}$.
\item Define $\mu:\mathcal{V}_{\mathbb{F}}\to K\otimes_{\mathbb{F}}\mathcal{V}_{\mathbb{F}}$ by $\mu(\mathbf{v})\defeq1\otimes\mathbf{v}$. Let $\mathcal{Y}$ be a $K$-vector space and let $f:\mathcal{V}_{\mathbb{F}}\to\mathcal{Y}$ be $\mathbb{F}$-linear.\\
A unique $K$-linear $\widetilde{f}:K\otimes_{\mathbb{F}}\mathcal{V}_{\mathbb{F}}\to\mathcal{Y}$ satisfies $\widetilde{f}\circ\mu=f$.
\end{enumerate}
\begin{proof}
\begin{enumerate}[label=(\roman*)]
\setlength{\itemsep}{0pt}
\item For $\alpha\in K$, set $m_\alpha(\beta,\mathbf{v})\defeq(\alpha\beta)\otimes\mathbf{v}$. The map $m_\alpha$ is $\mathbb{F}$-bilinear. It induces $M_\alpha\in\LL_{\mathbb{F}}(K\otimes_{\mathbb{F}}\mathcal{V}_{\mathbb{F}})$ by Theorem \ref{thm:tensor-linearizes-bilinear-maps}(ii).\\
For $\alpha,\lambda\in K$, the identities $M_0=0$ and $M_{\alpha+\lambda}=M_\alpha+M_\lambda$ hold on decomposable tensors.\\
Identities $M_{\alpha\lambda}=M_\alpha M_\lambda$ and $M_1=I$ likewise hold there. They extend to tensor product by Theorem \ref{thm:tensor-product-quotient-construction}(i).\\
The map $\alpha\mapsto M_\alpha$ is a unital ring homomorphism $K\to\operatorname{End}_{\mathbb{F}}(K\otimes_{\mathbb{F}}\mathcal{V}_{\mathbb{F}})$ by these identities.\\
The action $\alpha\cdot\mathbf{z}\defeq M_\alpha\mathbf{z}$ defines the claimed $K$-vector space by this unital ring homomorphism.
\item The map $\mu$ is $\mathbb{F}$-linear by bilinearity of the tensor map. Set $g(\beta,\mathbf{v})\defeq\beta f(\mathbf{v})$.\\
The map $g$ is $\mathbb{F}$-bilinear by the $\mathbb{F}$-linearity of $f$ and scalar multiplication on $\mathcal{Y}$.\\
A unique $\mathbb{F}$-linear map $\widetilde{f}$ satisfies $\widetilde{f}(\beta\otimes\mathbf{v})=\beta f(\mathbf{v})$ by Theorem \ref{thm:tensor-linearizes-bilinear-maps}(ii).\\
For $\alpha,\beta\in K$ and $\mathbf{v}\in\mathcal{V}_{\mathbb{F}}$, $\widetilde{f}(\alpha\cdot(\beta\otimes\mathbf{v}))=\alpha\beta f(\mathbf{v})=\alpha\widetilde{f}(\beta\otimes\mathbf{v})$.\\
The map $\widetilde{f}$ is $K$-linear by Theorem \ref{thm:tensor-product-quotient-construction}(i). The identity $\widetilde{f}\circ\mu=f$ holds by $\widetilde{f}(1\otimes\mathbf{v})=f(\mathbf{v})$.\\
If $h$ is $K$-linear, $h\circ\mu=f$, then $h(\beta\otimes\mathbf{v})=\beta h(1\otimes\mathbf{v})=\beta f(\mathbf{v})=\widetilde{f}(\beta\otimes\mathbf{v})$. Equality $h=\widetilde{f}$ holds by Theorem \ref{thm:tensor-product-quotient-construction}(i).
\end{enumerate}
\hlt{Structural Commentary}: Field-action maps define $K$-scalars; bilinear extension proves universality.\\
\hlt{Key Insight}: The factor $K$ freely supplies the missing scalars.
\end{proof}
\end{proposition}
